\section*{الفصل \ref{ch:g8:ohms-law} --- المقاومة وقانون أوم}

\begin{solution}{exo:g8:ohms-law:1}
مدى شدّة \omterm{def:g8:ohms-law:resistance}{مقاومة} العنصر للتيار؛ ووحدتها \omterm{def:g8:ohms-law:resistance}{الأوم}، ورمزها $\Omega$.
و$U = R I$؛ \ و$I = U/R$؛ \ و$R = U/I$.
\end{solution}

\begin{solution}{exo:g8:ohms-law:2}
$I = 6 \div 30 = \qty{0.2}{A}$؛ \ و$U = 30 \times 0.5 =
\qty{15}{V}$.
\end{solution}

\begin{solution}{exo:g8:ohms-law:3}
$R = 4.5 \div 0.09 = 50\,\Omega$.
\end{solution}

\begin{solution}{exo:g8:ohms-law:4}
خطّ مستقيم يمرّ بالمبدأ --- التناسب. والخطّ الأشدّ ميلًا يعني
فولتات أكثر لكل \omterm{def:g8:intensity-ammeter:intensity}{أمبير}: أي \omterm{def:g8:ohms-law:resistance}{مقاومة} أكبر.
\end{solution}

\begin{solution}{exo:g8:ohms-law:5}
$3.0 \div 0.20 = 15$؛ و$6.0 \div 0.40 = 15$. وتساوي النسبة عبر
السطور \emph{هو} القانون: فوجود $R$ ثابتة واحدة تخدم كل قراءة
هو ما يجعل ``\omterm{def:g8:ohms-law:resistance}{المقاومة}'' خاصّية من خواصّ العنصر.
\end{solution}

\begin{solution}{exo:g8:ohms-law:6}
\omterm{def:g8:ohms-law:resistance}{المقاومة} الأصغر تشرب أكثر: $230 \div 25 = \qty{9.2}{A}$ في
مقابل $230 \div 35 \approx \qty{6.6}{A}$.
\end{solution}

\begin{solution}{exo:g8:ohms-law:7}
المصباحان \omterm{def:g4:series-circuits:series}{على التسلسل} يتقاسمان $I$ واحدًا؛ \omterm{def:g8:voltage-voltmeter:voltage}{وتوتر} كل \omterm{def:g3:first-electric-circuit:bulb}{مصباح}
$U = R I$ بالتيار $I$ نفسه --- فتضرب $R$ الأكبر في العدد نفسه فتعطي
$U$ الأكبر: \omterm{def:g8:ohms-law:resistor}{فالمقاوم} الأقوى يأخذ الحصة الأكبر.
\end{solution}

\begin{solution}{exo:g8:ohms-law:8}
\omterm{def:g6:circuits-and-safety:short}{الدارة القصيرة}: $R$ قرب الصفر، فتنفجر $I = U/R$ --- فالاندفاع
قسمة على ما يكاد يكون لا شيء. \omterm{ex:g3:first-electric-circuit:switch}{والقاطعة} المفتوحة: $R$ فوق كل
\omterm{def:g6:measurement-in-science:measuring}{قياس}، فتنهار $I = U/R$ إلى الصفر --- الحصار التامّ.
\end{solution}

\begin{solution}{exo:g8:ohms-law:9}
يجب أن يأخذ \omterm{def:g8:ohms-law:resistor}{المقاوم} $9 - 2 = \qty{7}{V}$ عند \qty{0.02}{A}:
$R = 7 \div 0.02 = 350\,\Omega$.
\end{solution}

\begin{solution}{exo:g8:ohms-law:10}
النقطة الأولى: $R = 1 \div 0.25 = 4\,\Omega$؛ والثانية:
$4 \div 0.5 = 8\,\Omega$ --- فلا قيمة واحدة تلائم. والقصة:
التيار النامي يسخّن الخيط، والمعدن الساخن يقاوم أكثر؛
فتنحني الصورة لأن العنصر يتغيّر وهو يعمل.
\end{solution}

\begin{solution}{exo:g8:ohms-law:11}
حلقة واحدة، و$I$ واحد: فالهبوطان يجتمعان في دفعة \omterm{def:g3:first-electric-circuit:battery}{البطارية}،
$6 = 10 I + 20 I = 30 I$، ومنه $I = \qty{0.2}{A}$. ثم
$U_{10} = 10 \times 0.2 = \qty{2}{V}$ و$U_{20} = 20 \times 0.2
= \qty{4}{V}$ --- ومجموعهما \qty{6}{V}، كما يقتضي قانون
الحلقة.
\end{solution}

\begin{solution}{exo:g8:ohms-law:12}
كل \omterm{def:g5:series-parallel-circuits:parallel}{فرع} يمتدّ على \qty{6}{V}: $I_{10} = 6 \div 10 =
\qty{0.6}{A}$، و$I_{20} = 6 \div 20 = \qty{0.3}{A}$؛ والمجموع
\qty{0.9}{A}. والبديل الواحد: $R = 6 \div 0.9 \approx
6.7\,\Omega$ --- أي \emph{أقلّ} من أيّ عضو. وبلغة الطرق:
طريقان بين البلدتين نفسيهما يحملان حركة أكثر مما يحمله أحدهما
وحده؛ فكل طريق متفرّع مضاف يسهّل العبور، فتهبط \omterm{def:g8:ohms-law:resistance}{مقاومة} الفريق
دون \omterm{def:g8:ohms-law:resistance}{مقاومة} أضعف أعضائه.
\end{solution}

\begin{solution}{pb:g8:ohms-law:1}
\textbf{1.} المغذّي \omterm{def:g8:intensity-ammeter:ammeter}{ومقياس التيار} والعنصر في حلقة تسلسل واحدة
--- نقطة المكس في الطريق؛ \omterm{def:g8:voltage-voltmeter:voltmeter}{ومقياس التوتر} على طرفَي العنصر ---
المسّاح جانبًا.
\textbf{2.} $R = 4.5 \div 0.15 = 30\,\Omega$.
\textbf{3.} خمسة في المئة من $33$ نحو $1.65$: فالمقبول من
$31.35$ إلى $34.65\,\Omega$. و$30\,\Omega$ للعيّنة الأولى تقع
خارجه: مرفوضة.
\textbf{4.} $I = 4.5 \div 150 = \qty{0.03}{A} =
\qty{30.0}{mA}$.
\textbf{5.} كل سطر يعيد $R = U/I = 30\,\Omega$ ($1.5 \div
0.05$، و$3.0 \div 0.10$، \dots): نسبة ثابتة، وصورة على خطّ
مستقيم --- فقيمة واحدة تلائم الجميع: $30\,\Omega$.
\textbf{6.} النسب تجري $5$ و$6.7$ و$8.3$ و$10\,\Omega$:
تتسلّق مع كل خطوة --- \omterm{def:g8:ohms-law:resistance}{فالمقاومة} تنمو كلما نما التيار؛ ولا وجود
\omterm{def:g8:ohms-law:resistance}{لمقاومة} $R$ واحدة.
\textbf{7.} كلما نما التيار سخن الخيط نحو التوهّج، والمعدن
الساخن يقاوم أكثر --- فتنحني الصورة صعودًا. وهكذا لا يقبل
اختبار الخطّ المستقيم في المنضدة إلا العناصر التي تستقرّ
$R$ فيها: فهو، بحكم البناء، كاشف للمقاومات.
\textbf{8.} بند ثبات \omterm{def:g3:temperature-thermometers:temperature}{درجة الحرارة}: فقانون \omterm{def:g8:ohms-law:resistance}{أوم} يعد بالتناسب
\emph{عند \omterm{def:g3:temperature-thermometers:temperature}{درجة حرارة} ثابتة} --- وسخّن عنصرًا في منتصف الصورة
ينجرف حتى \omterm{def:g8:ohms-law:resistor}{المقاوم} الأمين.
\textbf{9.} $R = 12 \div 0.03 = 400\,\Omega$.
\textbf{10.} $U = R I = 25 \times 9.2 = \qty{230}{V}$ ---
أي الشبكة بالضبط: فالوشيعة تسحب تمامًا ما يخصّصه لها قانون
\omterm{def:g8:ohms-law:resistance}{أوم}، وهامش \omterm{def:g7:short-circuits-safety:fuse}{المصهر} يعمل كما صُمّم. فليس هنا ``لا'' ولا
تقصير؛ بل هو الحساب كله.
\textbf{11.} وُلدت الوصلة لتكون سلك ربط --- الطريق المستوي
الذي لا ينفق دفعة. وإن أُسيء استعمالها على طرفَي مغذٍّ، جعلتها
أوماتها شبه المعدومة دارةً قصيرة مثالية: فما كان الشرّير قطّ
إلا سلكًا جيّدًا جدًّا في المكان الخطأ.
\textbf{12.} مثلًا: ``القانون: $U = R I$، فالدفعة تساوي
\omterm{def:g8:ohms-law:resistance}{المقاومة} في المسير. وصورته: خطّ مستقيم يمرّ بالمبدأ، أشدّ
ميلًا كلما قوي \omterm{def:g8:ohms-law:resistor}{المقاوم}. ولا يجوز أن يوقّعه إلا عنصر عند \omterm{def:g3:temperature-thermometers:temperature}{درجة
حرارة} ثابتة --- \omterm{def:g5:heat-and-insulation:heat}{فالحرارة} تعيد كتابة \omterm{def:g8:ohms-law:resistance}{المقاومة}، والصور المأخوذة
في الحمّى لا تثبت شيئًا.''
\end{solution}
